\subsection{Vertical orders and horizontal section counts}
\label{subsec:rank:count}

We now apply the low-order families to the divisor supplied by
Lemma~\ref{lem:rank:calibration}.  Keep the contraction $h:X\to Y$, its
bounded comparison model $\widehat Y$, and the very ample divisor $H$ on
$\widehat Y$.  Write $b=\dim Y$, $s=n-b$, $A=h^*H$, and $N=H^b$.
Thus $0<b<n$, and both $N$ and
$|K_{\widehat Y}\cdot H^{b-1}|$ are bounded in terms of $n$.
There is no divisor of $X$ mapping into the codimension-two sets omitted
when identifying $Y$ with $\widehat Y$.
Let $D$ and $\delta>0$ be the calibrated divisor and parameter.  In this
subsection $C$ denotes constants depending only on $n$, independently of
$0<\eta\le1$; the threshold for a sufficiently large divisible integer
$m$ may depend on all the fixed geometric data.

For sufficiently large divisible $m$, our target is the estimate
\begin{equation}\label{eq:rank:count-target}
 h^0(X,\OO_X(mD))
 \le C\bigl(m\eta v(D)\bigr)^s(m\delta)^b.
\end{equation}
The first factor will bound, over the base function field, the number
of independent restrictions to the generic fibre, using vertical jets.
The second will bound their
contribution to global sections, using slopes on the base and a section
estimate linear in the sheaf rank.  Since $v(D)$ is comparable to
$\delta$, this would contradict the volume asymptotic
$h^0(X,mD)\sim m^n v(D)^n/n!$ when $\eta$ is sufficiently small.

\paragraph{The families are vertical.}
Apply Lemma~\ref{lem:rank:low-order} to the ample model $X_D$ of $D$.
If a geometric generic member $C_t$ has nonconstant image in
$\widehat Y$, choose a hyperplane in $|H|$ through the image of its marked
point $x_t$ without containing the image of $C_t$.  Since
$D-\delta A/3$ is big, choose a section of a sufficiently divisible
multiple of this divisor whose divisor avoids the generic mark.  This
section,
multiplied by the corresponding power of the hyperplane section, gives
\[
  \gamma(P_D|_{C_t};C_t,x_t)\ge \delta/3.
\]
Here one works on a common resolution and uses
\eqref{eq:rank:ample-model}.  The mark lies in the birational isomorphism
locus and avoids the fixed divisor; the chosen section remains nonzero
on the member.  Hyperplane pullbacks extend in codimension one on $X$
because no divisor maps into the omitted base sets.  The same
construction is valid over the geometric generic coefficient field.
The low-order bound is at most $\eta v(D)\le C\eta\delta$, contradicting
this inequality when $\eta$ is sufficiently small in terms of $n$.

The family is therefore vertical for $h$.  Its chain field contains
$\C(Y)$ and is attained.  Both fields are relatively algebraically
closed in $\C(X)$; maximality of $\dim Y$ among proper attained fields
forces them to coincide.  Indeed the chain field is proper by
Input~\ref{input:rank:chains} and the choice $\eta<1/n$, and an inclusion
of equal transcendence degree would be algebraic.  On a smooth model
$F$ of the geometric generic fibre of $h$, also resolving the ample-model
map, the chain estimate \eqref{eq:rank:chain-bound} now gives
\begin{equation}\label{eq:rank:vertical-order}
  \gamma(P_D|_F;F)\le s\eta v(D).
\end{equation}

For sufficiently divisible $m$, put $V_m=H^0(X,\OO_X(mD))$ and let $r_m$
be the dimension over $\C(Y)$ of its span in the generic fibre of
$h_*\OO_X(mD)$.  After extending to the geometric generic field, this
span identifies with a subspace of $H^0(F,mP_D|_F)$ multiplied by the
fixed part in \eqref{eq:rank:ample-model}.  At a general mark that fixed
part is a unit.  Thus \eqref{eq:rank:vertical-order} bounds the order of
every nonzero element of the span, including arbitrary scalar
combinations over that field.  Its vertical $k$-jet map is injective for
\[
  k=\lceil ms\eta v(D)\rceil.
\]
Consequently, for $m$ sufficiently large,
\begin{equation}\label{eq:rank:vertical-rank}
  r_m\le {s+k\choose s}\le C\bigl(m\eta v(D)\bigr)^s.
\end{equation}
Saturate the subsheaf generically spanned by $V_m$ in $h_*\OO_X(mD)$,
transfer it across the common big open of $Y$ and $\widehat Y$, and extend
reflexively.  This gives a reflexive sheaf $\mathcal E_m$ on
$\widehat Y$ of rank $r_m$, with
\begin{equation}\label{eq:rank:sections-in-span}
  \dim V_m\le h^0(\widehat Y,\mathcal E_m).
\end{equation}
We must bound the number of these sections without losing a further
power of $r_m$.

\paragraph{A slope bound from jets along the fibres.}
For a torsion-free sheaf $\mathcal S$ on $\widehat Y$ write
\[
  \mu_H(\mathcal S)
  =\frac{c_1(\mathcal S)\cdot H^{b-1}}{\rank\mathcal S}.
\]
Let $\mathcal S\subset\mathcal E_m$ be saturated of rank $d>0$.
On a smooth big open of $X$, the saturated relative tangent sheaf
$\mathcal T=\ker(dh)$ and its quotient in $\mathcal T_X$ are vector
bundles.  We remove their codimension-two defects, while retaining the
generic points of divisors where the differential of $h$ drops rank.
The evaluation map from $h^*\mathcal S$ to
$L=\OO_X(mD)$ is defined there.  Since $\mathcal T$ is integrable, the
bundles of jets along $\mathcal T$ have a filtration with quotients
\[
  L\otimes\Sym^i(\mathcal T^*),\qquad 0\le i\le k.
\]
For completeness, these bundles can be constructed by taking ordered
derivatives in local frames of $\mathcal T$.  The product rule makes
changes of frames triangular in the order; integrability identifies the
degree-$i$ symbol with $\Sym^i(\mathcal T^*)$.  The cocycle identities
agree with ordinary fibrewise jets on the dense smooth locus of $h$
and hence hold throughout this open.

Functions from the base have zero derivative along $\mathcal T$.
Taking jets of the evaluation map therefore gives an $\OO_X$-linear map
\[
  h^*\mathcal S\longrightarrow J^k_{\mathcal T}(L),
\]
generically injective by \eqref{eq:rank:vertical-order}.  Take its top
exterior power and a nonzero component in the induced graded
filtration.  In characteristic zero, the symmetric and exterior powers
occurring there embed in direct sums of tensor powers.  Selecting a
nonzero summand gives
\begin{equation}\label{eq:rank:jet-determinant}
  h^*\det\mathcal S\longrightarrow
       (\mathcal T^*)^{\otimes j}\otimes L^d,
  \qquad 0\le j\le kd.
\end{equation}

Wedge with the pulled-back base differentials defines a regular map
\begin{equation}\label{eq:rank:base-wedge}
  \mathcal T^*\longrightarrow
  \bigwedge^{b+1}\Omega_X^1\otimes\OO_X(-h^*K_{\widehat Y}),
\end{equation}
which is generically injective.  To check regularity, extend a local
leafwise covector to an ordinary covector and wedge it with the pulled-back
base volume form.  Two extensions differ by a conormal covector, whose
wedge is zero.  This also proves regularity along multiple-fibre
divisors: the pulled-back volume form may vanish there but has no pole.

Combining \eqref{eq:rank:jet-determinant} and
\eqref{eq:rank:base-wedge}, and moving the twists to the source, gives a
nonzero map from a line bundle of class
\begin{equation}\label{eq:rank:source-line}
  h^*\det\mathcal S-dmD+jh^*K_{\widehat Y}
\end{equation}
to a tensor power of $\Omega_X^1$.  We test this map by the functional
$\ell_X$ from Lemma~\ref{lem:rank:calibration}.  On a smooth resolution
$p:Z\to X$, define $\Lambda(M)=\ell_X(p_*M)$.  This is nonnegative on
pseudo-effective divisor classes and vanishes on exceptional classes.
The canonical class of $Z$ is effective and exceptional, since $X$ is
canonical and $K_X\sim0$, so $\Lambda(K_Z)=0$.

Saturate the line defined by the map on $Z$.  Campana--P\u{a}un's
cotangent theorem \cite[Theorem~1.3]{CP2019} says that every torsion-free
quotient of a cotangent tensor power on $Z$ has pseudo-effective
determinant.  The determinant of the tensor power itself is a multiple
of $K_Z$.  It follows that the saturated line has $\Lambda$-degree at
most zero.  Its pushforward differs from \eqref{eq:rank:source-line}
by an effective divisor, because the original map is regular in
codimension one.  Thus
\begin{align}
  \ell(\det\mathcal S)
   &\le dm\ell_X(D)-j\ell(K_{\widehat Y})
        \notag\\
   &\le Cdm\delta.\label{eq:rank:slope-bound}
\end{align}
The last inequality uses $j\le kd$, $\eta\le1$, the bounded canonical
degree of $\widehat Y$, and $v(D)\le C\delta$; the ceiling in $k$ is
absorbed once $m\delta\ge1$.  If $j=0$, the same conclusion follows
directly from the effective divisor of the map to $\OO_X$.
Since $\mathcal S$ was arbitrary, this proves
$\mu_{\max,H}(\mathcal E_m)\le Cm\delta$.

\paragraph{Counting sections without a rank loss.}
The following elementary consequence of the Grauert--M\"ulich--Spindler theorem
removes the rank from inside the polynomial section bound.

\begin{lemma}\label{lem:rank:section-count}
Let $b\ge1$.  A torsion-free slope-semistable sheaf $\mathcal G$ on
$\PP^b_{\C}$, of rank $u$ and slope $\mu$ with respect to
$\OO_{\PP^b}(1)$, satisfies
\[
  h^0(\PP^b,\mathcal G)
    \le C_bu\bigl(1+\max\{0,\mu\}\bigr)^b,
\]
where $C_b$ depends only on $b$.
\end{lemma}

\begin{proof}
If $\mu<0$, a nonzero section would yield a subsheaf of slope at least
zero, contradicting semistability.  Suppose $\mu\ge0$.
For $b\ge2$, the Grauert--M\"ulich--Spindler theorem
\cite{GrauertMulich1975,Spindler1979}, in its torsion-free form
\cite[Corollary~3.1.6]{HuybrechtsLehn2010}, bounds the highest degree $a$
of the splitting on a general line by $a\le\mu+(u-1)/2\le\mu+u$.
For $b=1$, the splitting theorem on $\PP^1$ and
semistability give the stronger equality $a=\mu$.

A general linear flag gives, for any torsion-free sheaf with highest
general-line splitting degree $a$,
\begin{equation}\label{eq:rank:preliminary-count}
  h^0(\mathcal G)\le u{a+b\choose b}\quad(a\ge0),
  \qquad h^0(\mathcal G)=0\quad(a<0).
\end{equation}
Indeed the successive restrictions can be chosen torsion-free: embed
the sheaf in a vector bundle and choose the cuts general for the
associated primes of the quotient as well.  Testing general lines
shows that twists by integers below $-a$ have no sections.  Starting
with the bound $u(a+1)$ on a line, the hyperplane restriction sequences
and the identity
$\sum_{i=0}^a {i+b-1\choose b-1}={a+b\choose b}$ prove
\eqref{eq:rank:preliminary-count} inductively.

For an integer $t\ge1$, let
\[
  f_t:\PP^b\longrightarrow\PP^b,
  \qquad [x_0:\cdots:x_b]\longmapsto[x_0^t:\cdots:x_b^t].
\]
Finite separable pullback preserves semistability
\cite[Lemma~2.1]{Weissmann2024}.  Since $f_t$ is flat,
$\mathcal G_t=f_t^*\mathcal G\otimes\OO_{\PP^b}(b(t-1))$ is
torsion-free and semistable, of rank $u$ and slope
$t\mu+b(t-1)$.  Moreover
\begin{equation}\label{eq:rank:power-sections}
  h^0(\mathcal G_t)\ge t^b h^0(\mathcal G).
\end{equation}
To see this, multiply the pulled-back sections by the $t^b$ homogeneous
monomials
\[
  x_0^{\,b(t-1)-\sum_{i=1}^b e_i}
  \prod_{i=1}^b x_i^{e_i},\qquad 0\le e_i<t.
\]
On $x_0\ne0$, these monomials form a basis of the upper function field
over the field generated by the $t$-th powers of the affine coordinates.
A linear relation among their products with pulled-back sections
therefore separates into linear relations among the original sections,
which proves \eqref{eq:rank:power-sections}.

Apply \eqref{eq:rank:preliminary-count} to $\mathcal G_t$, whose highest
general-line degree is at most $t\mu+b(t-1)+u$.  Divide the resulting
bound by $t^b$ and let $t\to\infty$, keeping $\mathcal G$ and $u$ fixed.
This yields
\[
  h^0(\mathcal G)\le \frac{u}{b!}(\mu+b)^b,
\]
and hence the asserted bound, for example with $C_b=b^b/b!$.
\end{proof}

\begin{proof}[Completion of the proof of Theorem~\ref{thm:rank}]
Choose a finite linear projection $f:\widehat Y\to\PP^b$ for the
embedding defined by $H$.  It has degree $N=H^b$ and
$f^*\OO_{\PP^b}(1)=\OO_{\widehat Y}(H)$.
The torsion-free sheaf $f_*\mathcal E_m$ has rank $Nr_m$ and
\begin{equation}\label{eq:rank:pushforward-slope}
  \mu_{\max}(f_*\mathcal E_m)\le Cm\delta/N.
\end{equation}
For this, let $\mathcal A$ be its first Harder--Narasimhan term.
Adjunction gives a nonzero map $f^*\mathcal A\to\mathcal E_m$.
After removing torsion and extending across codimension two, it gives
a map from the reflexive pullback $f^{[*]}\mathcal A$ to the reflexive
sheaf $\mathcal E_m$; denote its torsion-free image by $\mathcal Q$.
Finite separable reflexive pullback preserves semistability and
multiplies slopes for these polarizations by $N$
\cite[Lemma~2.1]{Weissmann2024}.
Consequently
\[
  N\mu(\mathcal A)
   \le\mu_H(\mathcal Q)
   \le\mu_{\max,H}(\mathcal E_m)
   \le Cm\delta,
\]
proving \eqref{eq:rank:pushforward-slope}.  No injectivity of the
adjunction map is required.

Apply Lemma~\ref{lem:rank:section-count} to the semistable
Harder--Narasimhan factors of $f_*\mathcal E_m$ and sum their section
bounds.  Their ranks sum to $Nr_m$, and $N$ is bounded in terms of $n$.
For all sufficiently large divisible $m$, this gives
\[
  \dim V_m\le h^0(\widehat Y,\mathcal E_m)
       \le Cr_m(m\delta)^b
       \le C\bigl(m\eta v(D)\bigr)^s(m\delta)^b.
\]
This is the target estimate \eqref{eq:rank:count-target}.
Divide by $m^n$ and use the volume asymptotic in divisible degrees.
Since $n=s+b$, we obtain
\[
  \frac{v(D)^n}{n!}\le C\eta^s v(D)^s\delta^b.
\]
The lower calibration bound $v(D)\ge c\delta$ then implies
$c^b/n!\le C\eta^s$.  This is impossible for a sufficiently small
positive $\eta$ depending only on $n$.  There are only finitely many
pairs $b,s>0$ with $b+s=n$, so choose $\eta$ for all of them at once,
also satisfying the earlier verticality and chain requirements, and
then impose the rank threshold in Lemma~\ref{lem:rank:low-order}.
The resulting contradiction bounds $r(T)$ in terms of $n$ alone.
\end{proof}
